\documentclass[12pt]{article}
\usepackage[letterpaper,margin=0.75in]{geometry}
\usepackage[T1]{fontenc}
\usepackage{lmodern}
\usepackage{microtype}
\usepackage{booktabs}
\usepackage{array}
\usepackage{amssymb}
\usepackage{enumitem}
\usepackage{hyperref}
\usepackage{setspace}
\usepackage{xcolor}
\setstretch{1.0}
\setlength{\parindent}{1em}
\setlength{\parskip}{0pt}
\setlength{\tabcolsep}{4pt}
\setlist{nosep,leftmargin=1.4em}
\hypersetup{
  hidelinks,
  pdftitle={One genre, two grammars: Community-conditioned subject omission in Git commit messages},
  pdfauthor={},
  pdfsubject={Anonymous PLC 51 abstract},
  pdfkeywords={}
}
\pagestyle{plain}
\renewcommand{\refname}{References}
\newcommand{\ex}[1]{\textit{#1}}

\begin{document}

\begin{center}
{\bfseries One genre, two grammars: Community-conditioned subject omission\\
in Git commit messages}
\end{center}
\vspace{-0.65em}

English permits null subjects in restricted registers, most prominently diaries
(Haegeman 1997, 2013; Weir 2012). Git commit subjects appear to add
a technologically mediated case: \ex{Fix parser} and \ex{Add tests} are routinely
called \emph{imperatives} in developer style guides. That label is not innocent.
An imperative analysis supplies both the missing subject (\textsc{you}) and the
base verb, whereas a register-drop analysis treats omission as independently
licensed by discourse and syntax. We show that the imperative diagnosis fails
as a general analysis. The same communicative task is realized by two nearly
categorical community conventions, one with a base verb and one with active
past tense. The contrast extends register-conditioned subject omission to
collaborative technical writing and shows that grammatical analysis of short
digital texts must include community norms and interface metadata.

\textbf{Data and method.}
We collected the latest 1,000 non-merge commit subjects from each of 14 major
public repositories (14,000 subjects total) on November 5, 2026. The sample
spans programming languages, web frameworks, and data-science projects.
Bracketed review tags and one repository scope prefix (e.g.,
\ex{net/http:}) were treated as metadata outside the clause. Automatic merges,
backports, and explicit reverts were excluded. To avoid relying on a general
parser trained on edited prose, we used a conservative, auditable
criterion: the first clausal token had to instantiate one of 40 hand-specified
English paradigms whose spelling distinguishes \textsc{base}, \textsc{past},
\textsc{3sg}, and \textsc{progressive} (e.g., \ex{fix/fixed/fixes/fixing}).
This retained 8,291 subjects (59.2\%). The criterion is transparent and exactly
reproducible from the public commit histories.

\textbf{Results.}
Of the retained subjects, 7,484 began with a base form, 752 with a past form,
23 with third-person singular, and 32 with a progressive. Table~\ref{tab:repos}
reports the decisive base--past comparison. Thirteen repositories selected the
base form in 97.4--100\% of tokens; Django selected the past in all 685
base--past tokens. The association between repository and form is extremely
large, $\chi^2(13)=7442.9$, $p<.001$, Cramer's $V=.951$ ($N=8{,}236$).
Lexical choice cannot explain the split: \textit{fix}, for example, occurs 2,213
times as \ex{fix} and 537 times as \ex{fixed} across the corpus. Nor are the
past forms reduced passives: attested strings such as (1b) contain an overt
direct object.

\begin{enumerate}[label=(\arabic*),leftmargin=2.2em,itemsep=0pt,topsep=2pt]
\item \begin{tabular}[t]{@{}ll@{}}
a. & \ex{Fix host mismatch in the test.}\\
b. & \ex{Fixed \#37388 -- Improved unsafe user input examples.}
\end{tabular}
\end{enumerate}
\vspace{-0.4em}

\begin{table}[h]
\centering
\caption{Clause-initial base and past forms by repository.}
\label{tab:repos}
\vspace{-0.35em}
\begin{tabular}{@{}lrr@{\hspace{1.35em}}lrr@{}}
\toprule
Repository & Base & Past & Repository & Base & Past\\
\midrule
CPython  & 564 & 1  & Requests & 643 & 16\\
React    & 496 & 3  & NumPy    & 460 & 6\\
Django   & 0   & 685& pandas   & 281 & 2\\
Go       & 467 & 2  & FastAPI  & 879 & 0\\
Node.js  & 455 & 0  & Vue      & 887 & 2\\
PyTorch  & 433 & 3  & Svelte   & 683 & 4\\
Flask    & 639 & 12 & HTTPX    & 597 & 16\\
\bottomrule
\end{tabular}
\end{table}
\vspace{-0.65em}

The Django distribution is independently corroborated by its public contributor
guide, which explicitly requires past tense and illustrates \ex{Fixed Unicode
bug} (Django Software Foundation 2026); several other public guides prescribe the
imperative base. These prescriptions are not noise imposed on grammar.
They are evidence that a speech community conventionalizes the realization of
a genre-specific construction.

\newpage
\textbf{Analysis.}
The past pattern cannot be an imperative: English imperatives lack past-tense
morphology, and the transitive objects in (1b) exclude a passive parse. It
instead instantiates written-register subject omission. The omitted argument
is recoverable as the responsible developer or team, made continuously salient
by the author field attached to every commit. Its clause is a root and the
gap is leftmost once non-clausal interface labels are removed, matching core
restrictions on diary drop (Haegeman 2013). Crucially, however, the
base pattern is not thereby proven directive. A commit subject reports a
completed change; it does not ordinarily direct its reader to perform one.
We therefore analyze the two patterns as competing allostructions of a
\textsc{commit-report} construction:

\[
\left[\textsc{commit-report}\;\varnothing_{\mathrm{author/team}}\;
V_{\{\mathrm{base},\mathrm{past}\}}\;XP\right].
\]

The genre licenses a recoverable left-edge subject, while repository-level
convention selects the verbal realization. On this account, \emph{imperative}
is a useful style label but an unreliable morphosyntactic diagnosis: base
morphology, null subject, and directive force come apart.

This analysis makes three testable predictions. First, authors moving between
repositories should accommodate their verb form to the destination community.
Second, removing visible authorship metadata should reduce the acceptability of
past-form commit reports more than overt-subject paraphrases. Third, material
that is syntactically clause-internal, unlike scope labels, should block
left-edge omission in the past construction. More broadly, the result supports
a view of digital registers in which platform fields participate in
recoverability, but community practice determines which grammatical resources
realize the same discourse function.

\vspace{-0.35em}
\begin{thebibliography}{9}
\setlength{\itemsep}{0pt}
\setlength{\parskip}{0pt}

\bibitem[Django Software Foundation 2026]{django2026}
Django Software Foundation. 2026.
\newblock Committing code.
\newblock \url{https://docs.djangoproject.com/en/dev/internals/contributing/committing-code/}.

\bibitem[Haegeman 1997]{haegeman2008}
Haegeman, Liliane. 1997.
\newblock Register variation, truncation, and subject omission in English and French.
\newblock \textit{English Language and Linguistics} 1(2). 233--270.

\bibitem[Haegeman 2013]{haegeman2013}
Haegeman, Liliane. 2013.
\newblock The syntax of registers: Diary subject omission and the privilege of the root.
\newblock \textit{Lingua} 130. 88--110.

\bibitem[Weir 2012]{weir2012}
Weir, Andrew. 2012.
\newblock Left-edge deletion in English and subject omission in diaries.
\newblock \textit{English Language and Linguistics} 16. 105--129.

\end{thebibliography}

\end{document}
